\section{Complexes}
\label{section-complexes}

\noindent
Of course the notions of a chain complex and a cochain complex
are dual and you only have to read one of the two parts of
this section. So pick the one you like. (Actually, this doesn't
quite work right since the conventions on numbering things
are not adapted to an easy transition between chain and cochain
complexes.)

\medskip\noindent
A {\it chain complex $A_\bullet$} in an additive category $\mathcal{A}$
is a complex
$$
\ldots \to
A_{n + 1} \xrightarrow{d_{n + 1}}
A_n \xrightarrow{d_n}
A_{n - 1} \to
\ldots
$$
of $\mathcal{A}$. In other words, we are given an object $A_i$ of
$\mathcal{A}$ for all $i \in \mathbf{Z}$ and for
all $i \in \mathbf{Z}$ a morphism $d_i : A_i \to A_{i - 1}$ such that
$d_{i - 1} \circ d_i = 0$ for all $i$. A {\it morphism of chain
complexes $f : A_\bullet \to B_\bullet$} is given by a
family of morphisms $f_i : A_i \to B_i$ such that all
the diagrams
$$
\xymatrix{
A_i \ar[r]_{d_i} \ar[d]_{f_i} & A_{i - 1} \ar[d]^{f_{i - 1}} \\
B_i \ar[r]^{d_i} & B_{i - 1}
}
$$
commute. The {\it category of chain complexes of $\mathcal{A}$}
is denoted $\text{Ch}(\mathcal{A})$. The full subcategory consisting
of objects of the form
$$
\ldots \to A_2 \to A_1 \to A_0 \to 0 \to 0 \to \ldots
$$
is denoted $\text{Ch}_{\geq 0}(\mathcal{A})$.
In other words, a chain complex $A_\bullet$ belongs to
$\text{Ch}_{\geq 0}(\mathcal{A})$ if and only if
$A_i = 0$ for all $i < 0$.

\medskip\noindent
Given an additive category $\mathcal{A}$ we identify $\mathcal{A}$
with the full subcategory of $\text{Ch}(\mathcal{A})$ consisting
of chain complexes zero except in degree $0$ by the functor
$$
\mathcal{A} \longrightarrow \text{Ch}(\mathcal{A}),\quad
A \longmapsto (\ldots \to 0 \to A \to 0 \to \ldots)
$$
By abuse of notation we often denote the object on the right hand side
simply $A$. If we want to stress that we are viewing $A$ as a chain
complex we may sometimes use the notation $A[0]$, see
Section \ref{section-homotopy-shift}.

\medskip\noindent
A {\it homotopy $h$} between a pair of morphisms
of chain complexes $f, g : A_\bullet \to B_\bullet$
is a collection of morphisms $h_i : A_i \to B_{i + 1}$
such that we have
$$
f_i - g_i = d_{i + 1} \circ h_i + h_{i - 1} \circ d_i
$$
for all $i$.
Two morphisms $f, g : A_\bullet \to B_\bullet$ are
said to be {\it homotopic} if a homotopy between $f$
and $g$ exists.
Clearly, the notions of chain complex, morphism of
chain complexes, and homotopies between morphisms of chain complexes
make sense even in a preadditive category.

\begin{lemma}
\label{lemma-compose-homotopy}
\begin{slogan}
Hom functors of $\text{Ch}(\mathcal{A})$ respect the homotopy relation.
\end{slogan}
Let $\mathcal{A}$ be an additive category.
Let $f, g : B_\bullet \to C_\bullet$ be morphisms
of chain complexes. Suppose given morphisms of chain
complexes $a : A_\bullet \to B_\bullet$, and
$c : C_\bullet \to D_\bullet$.
If $\{h_i : B_i \to C_{i + 1}\}$ defines a homotopy
between $f$ and $g$, then $\{c_{i + 1} \circ h_i \circ a_i\}$
defines a homotopy between $c \circ f \circ a$ and
$c \circ g \circ a$.
\end{lemma}

\begin{proof}
Omitted.
\end{proof}

\noindent
In particular this means that it makes sense to define
the category of chain complexes with maps up to homotopy.
We'll return to this later.

\begin{definition}
\label{definition-homotopy-equivalent}
Let $\mathcal{A}$ be an additive category.
We say a morphism $a : A_\bullet \to B_\bullet$
is a {\it homotopy equivalence} if there exists
a morphism $b : B_\bullet \to A_\bullet$
such that there exists a homotopy between
$b \circ a$ and $\text{id}_A$
and there exists a homotopy between $a \circ b$ and $\text{id}_B$.
If there exists such a morphism between $A_\bullet$ and $B_\bullet$, then
we say that $A_\bullet$ and $B_\bullet$ are {\it homotopy equivalent}.
\end{definition}

\noindent
In other words, two complexes are homotopy equivalent if they become
isomorphic in the category of complexes up to homotopy.

\begin{lemma}
\label{lemma-cat-chain-abelian}
Let $\mathcal{A}$ be an abelian category.
\begin{enumerate}
\item The category of chain complexes in $\mathcal{A}$ is
abelian.
\item A morphism of complexes
$f : A_\bullet \to B_\bullet$ is injective
if and only if each $f_n : A_n \to B_n$ is injective.
\item A morphism of complexes
$f : A_\bullet \to B_\bullet$ is surjective
if and only if each $f_n : A_n \to B_n$ is surjective.
\item A sequence of chain complexes
$$
A_\bullet \xrightarrow{f} B_\bullet \xrightarrow{g} C_\bullet
$$
is exact at $B_\bullet$ if and only if each sequence
$$
A_i \xrightarrow{f_i} B_i \xrightarrow{g_i} C_i
$$
is exact at $B_i$.
\end{enumerate}
\end{lemma}

\begin{proof}
Omitted.
\end{proof}

\noindent
For any $i \in \mathbf{Z}$ the $i$th {\it homology group}
of a chain complex $A_\bullet$ in an abelian category is defined by
the following formula
$$
H_i(A_\bullet) = \Ker(d_i)/\Im(d_{i + 1}).
$$
If $f : A_\bullet \to B_\bullet$ is a morphism of chain
complexes of $\mathcal{A}$ then we get an induced
morphism $H_i(f) : H_i(A_\bullet) \to H_i(B_\bullet)$
because clearly
$f_i(\Ker(d_i : A_i \to A_{i - 1})) \subset
\Ker(d_i : B_i \to B_{i - 1})$, and similarly
for $\Im(d_{i + 1})$.
Thus we obtain a functor
$$
H_i : \text{Ch}(\mathcal{A}) \longrightarrow \mathcal{A}.
$$

\begin{definition}
\label{definition-quasi-isomorphism}
Let $\mathcal{A}$ be an abelian category.
\begin{enumerate}
\item A morphism of chain complexes $f : A_\bullet \to B_\bullet$
is called a {\it quasi-isomorphism} if the induced
map $H_i(f) : H_i(A_\bullet) \to H_i(B_\bullet)$
is an isomorphism for all $i \in \mathbf{Z}$.
\item A chain complex $A_\bullet$ is called
{\it acyclic} if all of its homology objects
$H_i(A_\bullet)$ are zero.
\end{enumerate}
\end{definition}


\begin{lemma}
\label{lemma-map-homology-homotopy}
Let $\mathcal{A}$ be an abelian category.
\begin{enumerate}
\item If the maps $f, g : A_\bullet \to B_\bullet$ are
homotopic, then the induced maps $H_i(f)$ and $H_i(g)$
are equal.
\item If the map $f : A_\bullet \to B_\bullet$ is a homotopy
equivalence, then $f$ is a quasi-isomorphism.
\end{enumerate}
\end{lemma}

\begin{proof}
Omitted.
\end{proof}

\begin{lemma}
\label{lemma-long-exact-sequence-chain}
Let $\mathcal{A}$ be an abelian category.
Suppose that
$$
0 \to
A_\bullet \to
B_\bullet \to
C_\bullet \to
0
$$
is a short exact sequence of chain complexes of $\mathcal{A}$.
Then there is a canonical long exact homology sequence
$$
\xymatrix{
\ldots & \ldots & \ldots \ar[lld] \\
H_i(A_\bullet) \ar[r] & H_i(B_\bullet) \ar[r] & H_i(C_\bullet) \ar[lld] \\
H_{i - 1}(A_\bullet) \ar[r] &
H_{i - 1}(B_\bullet) \ar[r] &
H_{i - 1}(C_\bullet) \ar[lld] \\
\ldots & \ldots & \ldots \\
}
$$
\end{lemma}

\begin{proof}
Omitted. The maps come from the Snake Lemma \ref{lemma-snake}
applied to the diagrams
$$
\xymatrix{
&
A_i/\Im(d_{A, i + 1}) \ar[r] \ar[d]^{d_{A, i}} &
B_i/\Im(d_{B, i + 1}) \ar[r] \ar[d]^{d_{B, i}} &
C_i/\Im(d_{C, i + 1}) \ar[r] \ar[d]^{d_{C, i}} &
0 \\
0 \ar[r] &
\Ker(d_{A, i - 1}) \ar[r] &
\Ker(d_{B, i - 1}) \ar[r] &
\Ker(d_{C, i - 1}) &
}
$$
\end{proof}

\noindent
A {\it cochain complex $A^\bullet$} in an additive category $\mathcal{A}$
is a complex
$$
\ldots \to
A^{n - 1} \xrightarrow{d^{n - 1}}
A^n \xrightarrow{d^n}
A^{n + 1} \to
\ldots
$$
of $\mathcal{A}$. In other words, we are given an object $A^i$ of
$\mathcal{A}$ for all $i \in \mathbf{Z}$ and for
all $i \in \mathbf{Z}$ a morphism $d^i : A^i \to A^{i + 1}$ such that
$d^{i + 1} \circ d^i = 0$ for all $i$. A {\it morphism of cochain
complexes $f : A^\bullet \to B^\bullet$} is given by a
family of morphisms $f^i : A^i \to B^i$ such that all
the diagrams
$$
\xymatrix{
A^i \ar[r]_{d^i} \ar[d]_{f^i} & A^{i + 1} \ar[d]^{f^{i + 1}} \\
B^i \ar[r]^{d^i} & B^{i + 1}
}
$$
commute. The {\it category of cochain complexes of $\mathcal{A}$}
is denoted $\text{CoCh}(\mathcal{A})$. The full subcategory consisting
of objects of the form
$$
\ldots \to 0 \to 0 \to A^0 \to A^1 \to A^2 \to \ldots
$$
is denoted $\text{CoCh}_{\geq 0}(\mathcal{A})$.
In other words, a cochain complex $A^\bullet$ belongs to the subcategory
$\text{CoCh}_{\geq 0}(\mathcal{A})$ if and only if
$A^i = 0$ for all $i < 0$.

\medskip\noindent
Given an additive category $\mathcal{A}$ we identify $\mathcal{A}$
with the full subcategory of $\text{CoCh}(\mathcal{A})$ consisting
of cochain complexes zero except in degree $0$ by the functor
$$
\mathcal{A} \longrightarrow \text{CoCh}(\mathcal{A}),\quad
A \longmapsto (\ldots \to 0 \to A \to 0 \to \ldots)
$$
By abuse of notation we often denote the object on the right hand side
simply $A$. If we want to stress that we are viewing $A$ as a cochain
complex we may sometimes use the notation $A[0]$, see
Section \ref{section-homotopy-shift}.

\medskip\noindent
A {\it homotopy $h$} between a pair of morphisms
of cochain complexes $f, g : A^\bullet \to B^\bullet$
is a collection of morphisms $h^i : A^i \to B^{i - 1}$
such that we have
$$
f^i - g^i = d^{i - 1} \circ h^i + h^{i + 1} \circ d^i
$$
for all $i$.
Two morphisms $f, g : A^\bullet \to B^\bullet$ are
said to be {\it homotopic} if a homotopy between $f$
and $g$ exists.
Clearly, the notions of cochain complex, morphism of
cochain complexes, and homotopies between morphisms of cochain complexes
make sense even in a preadditive category.

\begin{lemma}
\label{lemma-compose-homotopy-cochain}
Let $\mathcal{A}$ be an additive category.
Let $f, g : B^\bullet \to C^\bullet$ be morphisms
of cochain complexes. Suppose given morphisms of cochain
complexes $a : A^\bullet \to B^\bullet$, and
$c : C^\bullet \to D^\bullet$.
If $\{h^i : B^i \to C^{i - 1}\}$ defines a homotopy
between $f$ and $g$, then $\{c^{i - 1} \circ h^i \circ a^i\}$
defines a homotopy between $c \circ f \circ a$ and
$c \circ g \circ a$.
\end{lemma}

\begin{proof}
Omitted.
\end{proof}

\noindent
In particular this means that it makes sense to define
the category of cochain complexes with maps up to homotopy.
We'll return to this later.

\begin{definition}
\label{definition-homotopy-equivalent-cochain}
Let $\mathcal{A}$ be an additive category.
We say a morphism $a : A^\bullet \to B^\bullet$
is a {\it homotopy equivalence} if there exists
a morphism $b : B^\bullet \to A^\bullet$
such that there exists a homotopy between
$b \circ a$ and $\text{id}_A$
and there exists a homotopy between $a \circ b$ and $\text{id}_B$.
If there exists such a morphism between $A^\bullet$ and $B^\bullet$, then
we say that $A^\bullet$ and $B^\bullet$ are {\it homotopy equivalent}.
\end{definition}

\noindent
In other words, two complexes are homotopy equivalent if they become
isomorphic in the category of complexes up to homotopy.

\begin{lemma}
\label{lemma-cat-cochain-abelian}
Let $\mathcal{A}$ be an abelian category.
\begin{enumerate}
\item The category of cochain complexes in $\mathcal{A}$ is
abelian.
\item A morphism of cochain complexes
$f : A^\bullet \to B^\bullet$ is injective
if and only if each $f^n : A^n \to B^n$ is injective.
\item A morphism of cochain complexes
$f : A^\bullet \to B^\bullet$ is surjective
if and only if each $f^n : A^n \to B^n$ is surjective.
\item A sequence of cochain complexes
$$
A^\bullet \xrightarrow{f} B^\bullet \xrightarrow{g} C^\bullet
$$
is exact at $B^\bullet$ if and only if each sequence
$$
A^i \xrightarrow{f^i} B^i \xrightarrow{g^i} C^i
$$
is exact at $B^i$.
\end{enumerate}
\end{lemma}

\begin{proof}
Omitted.
\end{proof}

\noindent
For any $i \in \mathbf{Z}$ the $i$th {\it cohomology group}
of a cochain complex $A^\bullet$ is defined by
the following formula
$$
H^i(A^\bullet) = \Ker(d^i)/\Im(d^{i - 1}).
$$
If $f : A^\bullet \to B^\bullet$ is a morphism of cochain
complexes of $\mathcal{A}$ then we get an induced
morphism $H^i(f) : H^i(A^\bullet) \to H^i(B^\bullet)$
because clearly
$f^i(\Ker(d^i : A^i \to A^{i + 1})) \subset
\Ker(d^i : B^i \to B^{i + 1})$, and similarly
for $\Im(d^{i - 1})$.
Thus we obtain a functor
$$
H^i : \text{CoCh}(\mathcal{A}) \longrightarrow \mathcal{A}.
$$

\begin{definition}
\label{definition-quasi-isomorphism-cochain}
Let $\mathcal{A}$ be an abelian category.
\begin{enumerate}
\item A morphism of cochain complexes $f : A^\bullet \to B^\bullet$
of $\mathcal{A}$ is called a {\it quasi-isomorphism} if the induced
maps $H^i(f) : H^i(A^\bullet) \to H^i(B^\bullet)$
is an isomorphism for all $i \in \mathbf{Z}$.
\item A cochain complex $A^\bullet$ is called
{\it acyclic} if all of its cohomology objects
$H^i(A^\bullet)$ are zero.
\end{enumerate}
\end{definition}

\begin{lemma}
\label{lemma-map-cohomology-homotopy-cochain}
Let $\mathcal{A}$ be an abelian category.
\begin{enumerate}
\item If the maps $f, g : A^\bullet \to B^\bullet$ are
homotopic, then the induced maps $H^i(f)$ and $H^i(g)$
are equal.
\item If $f : A^\bullet \to B^\bullet$ is a homotopy equivalence,
then $f$ is a quasi-isomorphism.
\end{enumerate}
\end{lemma}

\begin{proof}
Omitted.
\end{proof}

\begin{lemma}
\label{lemma-long-exact-sequence-cochain}
\begin{slogan}
Short exact sequences of complexes give rise to long exact sequences
of (co)homology.
\end{slogan}
Let $\mathcal{A}$ be an abelian category. Suppose that
$$
0 \to
A^\bullet \to
B^\bullet \to
C^\bullet \to
0
$$
is a short exact sequence of cochain complexes of $\mathcal{A}$.
Then there is a long exact cohomology sequence
$$
\xymatrix{
\ldots & \ldots & \ldots \ar[lld] \\
H^i(A^\bullet) \ar[r] &
H^i(B^\bullet) \ar[r] &
H^i(C^\bullet) \ar[lld] \\
H^{i + 1}(A^\bullet) \ar[r] &
H^{i + 1}(B^\bullet) \ar[r] &
H^{i + 1}(C^\bullet) \ar[lld] \\
\ldots & \ldots & \ldots \\
}
$$
The construction produces long exact cohomology sequences
which are functorial in the short exact
sequence and compatible with shifts as in
Definition \ref{definition-cohomology-shift}.
\end{lemma}

\begin{proof}
For the horizontal maps $H^i(A^\bullet) \to H^i(B^\bullet)$ and
$H^i(B^\bullet) \to H^i(C^\bullet)$ we use the fact that $H^i$
is a functor, see above. For the ``boundary map''
$H^i(C^\bullet) \to H^{i + 1}(A^\bullet)$ we use the map $\delta$
of the Snake Lemma \ref{lemma-snake}
applied to the diagram
$$
\xymatrix{
&
A^i/\Im(d_A^{i - 1}) \ar[r] \ar[d]^{d_A^i} &
B^i/\Im(d_B^{i - 1}) \ar[r] \ar[d]^{d_B^i} &
C^i/\Im(d_C^{i - 1}) \ar[r] \ar[d]^{d_C^i} &
0 \\
0 \ar[r] &
\Ker(d_A^{i + 1}) \ar[r] &
\Ker(d_B^{i + 1}) \ar[r] &
\Ker(d_C^{i + 1}) &
}
$$
This works as the kernel of the right vertical map is equal to
$H^i(C^\bullet)$ and the cokernel of the left vertical map is
$H^{i + 1}(A^\bullet)$. The exactness of the long sequence is
the exactnesss in part (2) of Lemma \ref{lemma-snake}.
The functoriality is Lemma \ref{lemma-snake-natural}. Compatibility
with shifts is immediate from the definitions.
\end{proof}








